% Partie : histoire
% Chapitre : biographies
% Section : Emmy Noether

\documentclass[../../../include/open-logic-section]{subfiles}

\begin{document}

\olfileid[fr]{his}{bio}{noe}

\olsection{Emmy Noether}

\olphoto{noether-emmy}{Emmy Noether}

Emmy Noether
(\textsc{neu}-teur)
naquit le
23 mars 1882
à Erlangen,
en Allemagne,
dans une famille
de lettrés
de la classe
moyenne supérieure.
Saluée comme
la «~mère
de l'algèbre
moderne~»,
elle apporta
des contributions
fondamentales
aux mathématiques
et à la physique
malgré les
obstacles
considérables
à l'éducation
des femmes.
Dans l'Allemagne
de l'époque,
les jeunes filles
devaient recevoir
une éducation
artistique et
ne pouvaient
pas entrer
dans les écoles
préparant aux
études universitaires.
Noether assista
cependant, en
auditrice libre,
à des cours
des universités
de G\"{o}ttingen
et d'Erlangen,
où son père
était professeur
de mathématiques.
En 1904,
après une
modification
de la politique
d'admission
d'Erlangen,
elle put enfin
s'y inscrire
comme étudiante.
Elle obtint
son doctorat
en mathématiques
en 1907.

Malgré ses
qualifications,
Noether se
heurta à de
fortes résistances
pendant sa
carrière.
De 1908 à 1915,
elle enseigna
à Erlangen
sans rémunération.
Son travail
attira alors
l'attention de
David Hilbert,
l'un des plus
grands mathématiciens
de l'époque,
qui l'invita
à G\"{o}ttingen.
Mais les femmes
n'y avaient pas
accès aux postes
de professeur.
Elle ne put
donner des cours
que sous le nom
de Hilbert,
là encore
sans salaire.
C'est durant
cette période
qu'elle démontra
ce qu'on appelle
aujourd'hui
le théorème
de Noether,
toujours utilisé
en physique
théorique.
Elle obtint
finalement le
droit d'enseigner
en 1919.
Hilbert aurait
répondu aux
objections persistantes
de ses collègues :
«~Messieurs,
le sénat de
la faculté
n'est pas
un établissement
de bains.~»

À la fin
des années 1920,
Noether se consacra
à l'algèbre
abstraite et
révolutionna
ce domaine.
Dans ses
démonstrations,
elle utilisait
souvent la
condition dite
de chaîne
ascendante :
pour certains
ensembles, il
n'existe pas
de chaîne
infinie strictement
croissante.
Ainsi, les
structures
algébriques
aujourd'hui
appelées anneaux
noethériens
n'admettent pas
de suite infinie
d'idéaux
$I_1 \subsetneq
I_2 \subsetneq
\dots$.
Cette condition
s'étend à
n'importe quel
ordre partiel
(en algèbre,
le cas particulier
est celui des
idéaux ordonnés
par inclusion).
On peut également
considérer sa
condition duale,
dite de chaîne
descendante :
toute suite
strictement
\emph{dé}croissante
dans l'ordre
partiel finit
par s'arrêter.
Lorsqu'un ordre
partiel satisfait
cette condition,
on peut raisonner
par induction
le long de
cet ordre,
comme on raisonne
par récurrence
selon l'ordre~$<$
sur~$\Nat$.
De tels ordres
sont dits
\emph{bien fondés}
ou \emph{noethériens},
et le principe
de démonstration
correspondant
s'appelle
\emph{induction
noethérienne}.

Noether était
juive. Lorsque
les nazis
prirent le
pouvoir en 1933,
elle fut
renvoyée
de son poste.
Elle put
heureusement
émigrer aux
États-Unis
pour occuper
un poste
temporaire
à Bryn Mawr,
en Pennsylvanie.
Elle donna
aussi des cours
à Princeton,
université
qu'elle jugea
peu accueillante
envers les femmes
\citep[81]{Dick1981}.
En 1935,
elle subit
une opération
pour retirer
une tumeur
de l'utérus.
Elle mourut
d'une infection
consécutive
à l'intervention
et fut
inhumée
à Bryn Mawr.

\begin{reading}
Pour une biographie
de Noether,
voir \citet{Dick1981}.
Le Perimeter
Institute for
Theoretical Physics
propose en ligne
des conférences
sur sa vie
et son influence
\citep{Perimeter2015}.
Si vous préférez
écouter plutôt
que lire,
\emph{Stuff You
Missed in History
Class} lui
consacre un
podcast
\citep{Frey2015}.
Ses œuvres
complètes sont
disponibles
dans l'original
allemand
\citep{Noether1983}.
\end{reading}

\end{document}
